\documentclass{article}
\usepackage{sagetex}
\title{A paper with its own helper code}
\author{Sagebrush}
\date{}
\begin{document}
\maketitle

\emph{In this preview:} this document is a folder: \texttt{paper.tex} and, beside it, \texttt{helper.py}, which
holds its longer code. The SageTeX blocks import it like any Python module, because a document's code runs in the
document's own folder, as \texttt{sage paper.sagetex.sage} does. Download or share it \emph{with its files}, and
the helper comes along.

\begin{sagesilent}
from helper import pell, collatz_steps
sols = pell(61, 2)
\end{sagesilent}

\section*{Pell's equation}
Fermat's challenge: the smallest solution of $x^2 - 61y^2 = 1$ in positive integers is
$(x, y) = \sage{sols[0]}$, and the next is $\sage{sols[1]}$.
Check: $x^2 - 61 y^2 = \sage{sols[1][0]^2 - 61*sols[1][1]^2}$.

\section*{Collatz}
Of the starting values below $1000$, the one taking the most steps to reach $1$ is
$n = \sage{max(range(1, 1000), key=collatz_steps)}$, after
$\sage{collatz_steps(max(range(1, 1000), key=collatz_steps))}$ steps.
\end{document}
